\section{Concluding remarks}\label{sec:conclusions}

This paper introduced \texttt{pyGLAM}, a Python package implementing the FKML and
Ramberg-Schmeiser parameterizations of the generalized lambda distribution behind
one interface, fitting them by multi-start moment matching and scoring the result
with an integrated goodness-of-fit module. A cross-validation against the R
packages and a factorial Monte Carlo study over two parameterizations, six
targets, four sample sizes and one hundred replicates per cell established four
things.

First, the implementation agrees with \texttt{gld} and \texttt{GLDEX}: the
quantile functions, distribution functions, densities and moments coincide to
numerical precision, and on the same samples the fits of \texttt{pyGLAM} are
valid, reproduce the sample moments, and take a fraction of the time. Where the
two packages reach different fits, the difference comes from how they solve the
moment equations and which root they reach, not from the functions they
evaluate.

Second, moment matching recovers all six targets to a practically useful
accuracy: at $N = 5{,}000$ the median error at $P_{50}$ is below $0.5\%$ of the
spread between $P_5$ and $P_{95}$ in ten of the twelve cases, at a mean cost of
well under a second per fit.

Third, the error separates into a sampling component and a structural one, and
the slope of the Kolmogorov-Smirnov distance against sample size tells them
apart. The normal and the uniform show no structural component under either
parameterization, the Gumbel and the Beta a small one, and the Gamma a floor
under both. On the Gamma the two parameterizations match the same moments with
distributions that fail in different places, so which one is better depends on
the metric that scores them; on the Gumbel the RS fits approach the exponential
tail through a degenerate limit of their parameters that the FKML form contains
as a regular case. On the Student-$t$ with four degrees of freedom both forms
slow down because the kurtosis of the target is infinite, which makes the
decomposition a diagnostic of the estimator as well. Since the Python
implementations we are aware of commit to a single parameterization, a
practitioner using them accepts whatever structural component that form carries
on their target, without being told. This is the argument for carrying more than
one form, and a diagnostic, in the same package.

Fourth, four-moment matching admits roots that reproduce the four sample moments,
are reported as converged by the optimizer, and are not distributions. Imposing
the admissible regions inside the RS objective, starting the multi-start on both
sides of those regions, and screening the candidates of both forms for a
non-decreasing quantile function kept every one of the $4{,}800$ selected fits
valid.

Several limitations point to the next steps. The package currently offers one
estimator, and moment matching is known to be the weakest of the GLD estimators;
the Student-$t$ case shows its limits directly. L-moments
\parencite{KarvanenNuutinen2008}, which exist whenever the mean does, percentile
matching \parencite{KarianDudewicz1999} and discretised maximum likelihood
\parencite{Su2007CSDA} are the natural additions, and the regime decomposition
used here transfers directly to comparing them. The comparison with
\texttt{GLDEX} shows that a denser search of starting points can reach better
roots of the RS moment equations on some targets, at a much higher cost; an
adaptive search that spends more starts only where the candidates disagree would
combine the two. The study covers six targets, chosen for geometric contrast;
multimodal targets, which a single GLD cannot represent, call for mixtures. The
\texttt{cdf} and \texttt{pdf} methods invert the quantile function elementwise
and would benefit from a vectorised inversion.
